\par\noindent\textbf{(1)} 对沿+z推动，横向变换E'\_⊥=γ(E\_⊥+u×B)，B'\_⊥=γ(B\_⊥−u×E/c²)，若采用S'相对S的标准变换这里u向量应取为−u ẑ；明确写坐标分量为E'\_x=γ(E−uB)、B'\_y=γ(B−uE/c²)，E'\_z=B'\_z=0。该符号对应S'与S相对速度+u ẑ。
\par\noindent\textbf{(2)} 直接展开E'\textasciicircum{}2−c²B'\textasciicircum{}2=γ²[(E−uB)²−c²(B−uE/c²)²]=γ²(1−u²/c²)(E²−c²B²)=E²−c²B²；仍有E'·B'=0，因为两场正交。该两个组合为Lorentz标量，而单独E大小和B大小不是。
\par\noindent\textbf{(3)} 若E<cB，可取u=E/B<c，使E'\_x=γ(E−uB)=0，剩余纯磁场B'\_y=√(B²−E²/c²)，由不变量定幅度。该推动沿E×B方向，表明磁占优正交场可以找到电场为零的惯性系。
\par\noindent\textbf{(4)} 若E>cB，令B'\_y=0得u=c²B/E<c，剩余电场E'\_x=√(E²−c²B²)，属于电占优场。当E=cB为零场不变量（null field），所需u→c，任何有限惯性推动都不能同时消掉电磁场，光波属于这种类型。
